\section*{अध्याय \ref{ch:g11:infrared} --- अवरक्त स्पेक्ट्रमिकी}

\begin{solution}{exo:g11:infrared:1}
$\lambda = 1/3300 = \qty{3.03e-4}{cm} = \qty{3.03}{\micro m}$;
$\lambda = 1/1050 = \qty{9.52e-4}{cm} = \qty{9.52}{\micro m}$। बड़ी \omterm{def:g11:infrared:wavenumber}{तरंग संख्या}
वाला \qty{3300}{cm^{-1}} का पट्ट तेज़ कंपन का है।
\end{solution}

\begin{solution}{exo:g11:infrared:2}
$\qty{10}{\micro m} = \qty{1.0e-3}{cm}$, इसलिए $\sigma = 1/\num{1.0e-3} =
\qty{1000}{cm^{-1}}$।
\end{solution}

\begin{solution}{exo:g11:infrared:3}
\omterm{def:g11:functional-groups:carbonyl}{कीटोन} का \ce{C=O}; हाइड्रोजन-आबंधित \omterm{def:g11:functional-groups:alcohol}{ऐल्कोहॉल} का \ce{O-H}; कार्बन श्रृंखला के
\ce{C-H} आबंध।
\end{solution}

\begin{solution}{exo:g11:infrared:4}
\qty{1715}{cm^{-1}} का पट्ट: \ce{C=O} आबंध।
\end{solution}

\begin{solution}{exo:g11:infrared:5}
उपकरण नमूने से पार निकलने वाले प्रकाश को मापता है; जहाँ कोई
आबंध अवशोषित करता है, वहाँ कम प्रकाश पार निकलता है और वक्र नीचे गिरता है। $T =
\qty{100}{\percent}$ का अर्थ है कि उस \omterm{def:g11:infrared:wavenumber}{तरंग संख्या} का कोई भी प्रकाश
अवशोषित नहीं होता।
\end{solution}

\begin{solution}{exo:g11:infrared:6}
पहले में \qty{1715}{cm^{-1}} पर एक \ce{C=O} पट्ट है और कोई \ce{O-H}
पट्ट नहीं। वह एक \omterm{def:g11:functional-groups:carbonyl}{कीटोन} है: \ce{CH3-CO-CH2-CH3}, ब्यूटेनोन (जिसे
ब्यूटेन-2-ओन भी कहते हैं)। दूसरे में एक \ce{O-H} पट्ट है और कोई \ce{C=O} नहीं: एक
\omterm{def:g11:functional-groups:alcohol}{ऐल्कोहॉल}, उदाहरण के लिए ब्यूटेन-1-ऑल।
\end{solution}

\begin{solution}{exo:g11:infrared:7}
अम्ल के \ce{O-H} समूह बहुत प्रबलता से हाइड्रोजन-आबंधित होते हैं (\omterm{def:g7:atoms-and-molecules:molecule}{अणु}
जोड़े बनाते हैं, हर \ce{O-H} दूसरे के \ce{C=O} से आबंधित),
और हर \omterm{def:g7:atoms-and-molecules:molecule}{अणु} थोड़ा अलग ढंग से आबंधित होता है: पट्ट और अधिक ढीला होता है
और एक बहुत चौड़े परास में फैल जाता है।
\end{solution}

\begin{solution}{exo:g11:infrared:8}
द्रव (A) में \ce{O-H} पट्ट चौड़ा है, लगभग
\qty{3350}{cm^{-1}} पर; गैस (B) में, जहाँ \omterm{def:g7:atoms-and-molecules:molecule}{अणु} एक-दूसरे से दूर होते हैं
और कोई \omterm{def:g11:polarity-and-cohesion:hydrogen-bond}{हाइड्रोजन आबंध} नहीं बनाते, वह लगभग
\qty{3666}{cm^{-1}} पर एक तीक्ष्ण पट्ट है।
\end{solution}

\begin{solution}{exo:g11:infrared:9}
मुश्किल से: \omterm{def:g11:functional-groups:carbonyl}{ऐल्डिहाइड} के लिए लगभग \qty{1730}{cm^{-1}} और
\omterm{def:g11:functional-groups:carbonyl}{कीटोन} के लिए \qty{1715}{cm^{-1}}, निकट मान। \omterm{def:g11:infrared:band}{अंगुलिछाप क्षेत्र} सहित पूरे
स्पेक्ट्रम की संदर्भ स्पेक्ट्रमों से तुलना
फ़ैसला करेगी, या कोई दूसरी तकनीक।
\end{solution}

\begin{solution}{exo:g11:infrared:10}
एक \omterm{def:g11:functional-groups:carboxylic-acid}{एस्टर} (\ce{C=O} लगभग \qty{1735}{cm^{-1}} पर, प्रबल \ce{C-O}, कोई
\ce{O-H} नहीं), उदाहरण के लिए एथिल एथेनोएट। एथेनोइक अम्ल \ce{C2H4O2} है,
\ce{C4H8O2} नहीं। ब्यूटेनोइक अम्ल का सूत्र सही है, पर वह अम्लों का बहुत
चौड़ा \ce{O-H} पट्ट दिखाता: बाहर।
\end{solution}

\begin{solution}{exo:g11:infrared:11}
प्रोपेन-1-ऑल एक \omterm{def:g11:functional-groups:alcohol}{ऐल्कोहॉल} है: A जैसा। प्रोपेनोइक अम्ल एक \omterm{def:g11:functional-groups:carboxylic-acid}{कार्बोक्सिलिक अम्ल} है:
D जैसा।
\end{solution}

\begin{solution}{exo:g11:infrared:12}
लगभग \qty{3350}{cm^{-1}} का चौड़ा \ce{O-H} पट्ट सिकुड़ता है, जबकि
लगभग \qty{1715}{cm^{-1}} का एक प्रबल \ce{C=O} पट्ट बढ़ता है। अभिक्रिया तब पूरी होती है
जब \ce{O-H} पट्ट गायब हो जाए और \ce{C=O} पट्ट और न
बढ़े।
\end{solution}

\begin{solution}{exo:g11:infrared:13}
ब्यूटेनोइक अम्ल 2500 से \qty{3300}{cm^{-1}} तक एक बहुत चौड़ा \ce{O-H} पट्ट
दिखाता है, \omterm{def:g11:functional-groups:carboxylic-acid}{एस्टर} कोई नहीं; अम्ल का \ce{C=O} लगभग
\qty{1710}{cm^{-1}} पर है, अधिक चौड़ा, \omterm{def:g11:functional-groups:carboxylic-acid}{एस्टर} का लगभग \qty{1735}{cm^{-1}} पर,
तीक्ष्ण, साथ में लगभग \qty{1240}{cm^{-1}} पर एक प्रबल \ce{C-O} पट्ट। अम्ल में
एक हाइड्रोजन-आबंधित \ce{O-H} समूह है, जो उसके \ce{C=O} को भी
थोड़ा ढीला कर देता है।
\end{solution}

\begin{solution}{exo:g11:infrared:14}
बचा हुआ एथेनॉल (वह \omterm{def:g11:functional-groups:alcohol}{ऐल्कोहॉल} जिससे \omterm{def:g11:functional-groups:carboxylic-acid}{एस्टर} बनता है) या पानी (\omterm{def:g4:air-a-mixture-of-gases:air}{वायु} से
या धुलाई से): दोनों में हाइड्रोजन-आबंधित \ce{O-H} समूह हैं।
नमूने की संदर्भ स्पेक्ट्रम से तुलना की जा सकती है, उसे सुखाया या आसवित किया जा सकता है, या
उसका क्वथन ताप मापा जा सकता है।
\end{solution}

\begin{solution}{exo:g11:infrared:15}
समान \omterm{def:g7:atoms-and-molecules:atom}{परमाणुओं} के बीच \omterm{def:g10:lewis-and-shape:multiple-bond}{द्विआबंध} एकल आबंध से अधिक कड़ी स्प्रिंग है:
वह तेज़ कंपन करता है, ऊँची \omterm{def:g11:infrared:wavenumber}{तरंग संख्या} पर (1715 के मुक़ाबले
\qty{1050}{cm^{-1}})। \ce{O-H}, \ce{N-H} और \ce{C-H} आबंधों सभी में
एक हाइड्रोजन \omterm{def:g7:atoms-and-molecules:atom}{परमाणु} है, जो सबसे हल्का है: स्प्रिंग पर हल्की गेंदें
तेज़ कंपन करती हैं, इसलिए लगभग 2900--\qty{3600}{cm^{-1}} की \omterm{def:g11:infrared:wavenumber}{तरंग संख्याएँ}।
\end{solution}

\begin{solution}{pb:g11:infrared:1}
\textbf{1.} \ce{CH3-CH2-OH}; \ce{CH3-CO-CH3}; \ce{CH3-COOH}।

\textbf{2.} \ce{C2H6O}; \ce{C3H6O}; \ce{C2H4O2}।

\textbf{3.} हाइड्रॉक्सिल; कार्बोनिल (\omterm{def:g11:functional-groups:carbonyl}{कीटोन}); कार्बोक्सिल।

\textbf{4.} एथेनॉल: \ce{O-H} और \ce{C-H}। प्रोपेनोन: \ce{C-H} और
\ce{C=O}। एथेनोइक अम्ल: \ce{O-H}, \ce{C-H} और \ce{C=O}।

\textbf{5.} कोई \ce{C=O} पट्ट नहीं, एक चौड़ा \ce{O-H} पट्ट: एथेनॉल।

\textbf{6.} 2500 से \qty{3300}{cm^{-1}} तक एक बहुत चौड़ा \ce{O-H} पट्ट
और लगभग \qty{1710}{cm^{-1}} पर एक प्रबल \ce{C=O}: एथेनोइक अम्ल।

\textbf{7.} प्रोपेनोन: \qty{1715}{cm^{-1}} पर बहुत प्रबल \ce{C=O} पट्ट,
बिना किसी \ce{O-H} के।

\textbf{8.} शायद (सिरका, \omterm{def:g6:solutions-and-solubility:solute}{विलायक}, \omterm{def:g11:functional-groups:alcohol}{ऐल्कोहॉल}), पर अज्ञात द्रवों को सूँघना
ख़तरनाक है; स्पेक्ट्रम सुरक्षित और निश्चित है।

\textbf{9.} अम्ल के \ce{O-H} समूह अधिक प्रबलता से हाइड्रोजन-आबंधित हैं
(जोड़े बनाए \omterm{def:g7:atoms-and-molecules:molecule}{अणु}), इसलिए पट्ट नीचे है और कहीं अधिक चौड़ा।

\textbf{10.} \ce{CH3-O-CH3}, \ce{C2H6O}; उसमें कोई \ce{O-H} नहीं, इसलिए कोई
\ce{O-H} पट्ट नहीं।

\textbf{11.} \ce{C=O} पट्ट: प्रोपेनैल के लिए लगभग \qty{1730}{cm^{-1}},
प्रोपेनोन के लिए लगभग \qty{1715}{cm^{-1}}।

\textbf{12.} \ce{O-H} पट्ट: \omterm{def:g11:functional-groups:carboxylic-acid}{एस्टर} में कोई \ce{O-H} नहीं होता। उसका \ce{C=O}
लगभग \qty{1735}{cm^{-1}} पर होगा।

\textbf{13.} \omterm{def:g11:organic-skeletons:isomer}{समावयवियों} का सूत्र एक ही होता है; उनके \omterm{def:g11:functional-groups:functional-group}{क्रियात्मक समूह} अलग होते हैं,
और स्पेक्ट्रम समूहों को दिखाता है।

\textbf{14.} \qty{1715}{cm^{-1}}।

\textbf{15.} $1715 \times 100 = \qty{1.715e5}{m^{-1}}$।

\textbf{16.} $\lambda = 1/\num{1.715e5} = \qty{5.83e-6}{m} =
\qty{5.83}{\micro m}$।

\textbf{17.} नहीं: दृश्य प्रकाश लगभग 0.4 से
\qty{0.75}{\micro m} तक जाता है। यह अवरक्त विकिरण है।

\textbf{18.} $\num{1.05e5}\ \si{m^{-1}}$, इसलिए $\lambda = \qty{9.52e-6}{m}
= \qty{9.52}{\micro m}$।

\textbf{19.} \ce{C=O} आबंध (ऊँची \omterm{def:g11:infrared:wavenumber}{तरंग संख्या}, छोटी तरंगदैर्घ्य)।

\textbf{20.} लगभग \qty{5.8}{\micro m}।
\end{solution}
